\subsection{Single Variable Calculus}

Calculus is the branch of mathematics that studies change and accumulation. It has two main branches:
\begin{enumerate}[label=\roman*.]
\setlength{\itemsep}{0pt}
\item Differential Calculus: Concerned with the rate at which things change (e.g. velocity).
\item Integral Calculus: Concerned with accumulation (e.g. area under a curve).
\end{enumerate}
Calculus underpins physics, engineering, economics, biology, and more.

\begin{definition} The \hlt{derivative} is defined as:
\begin{align*}
\lim_{\Delta x \rightarrow 0} \frac{\Delta f}{\Delta x} = \lim_{\Delta x \rightarrow 0} \frac{f(x_0 + \Delta x) - f(x_0)}{\Delta x} = f'(x_0) \nonumber
\end{align*}
\end{definition}

\begin{definition} A function $f(x)$ is continuous at $x_0$ when
\begin{align}
\lim_{x \rightarrow x_0} f(x) = f(x_0) \nonumber
\end{align}
\end{definition}

\begin{definition} The \hlt{right-hand} and \hlt{left-hand} limits are as follows:
\begin{enumerate}[label=\roman*.]
\setlength{\itemsep}{0pt}
\item \hlt{Right-Hand Limit}: $\lim\limits_{x \rightarrow x^+_0} f(x)$ means $\lim\limits_{x \rightarrow x_0} f(x)$ for $x > x_0$
\item \hlt{Left-Hand Limit}: $\lim\limits_{x \rightarrow x^-_0} f(x)$ means $\lim\limits_{x \rightarrow x_0} f(x)$ for $x < x_0$
\end{enumerate}
\end{definition}

\begin{definition} A \hlt{removable discontinuity} is where
\begin{align}
\lim\limits_{x \rightarrow x^+_0} f(x) = \lim\limits_{x \rightarrow x^-_0} f(x) \nonumber
\end{align}
but not at $f(x_0)$, or if $f(x_0)$ is undefined.
\end{definition}

\begin{definition} A \hlt{jump discontinuity} is where $\lim\limits_{x \rightarrow x^+_0} f(x)$ and $\lim\limits_{x \rightarrow x^-_0} f(x)$ exists, but not equal.
\end{definition}

\begin{definition} A \hlt{infinite discontinuity} is where $\lim\limits_{x \rightarrow x^+_0} f(x) \neq \lim\limits_{x \rightarrow x^-_0} f(x) \in \{+\infty, -\infty\}$
\end{definition}

\begin{theorem} \hlt{Differentiability Implies Continuity}\\
If $f$ is differentiable at $x_0$, then $f$ is continuous at $x_0$.
\begin{proof}
\begin{align}
\lim_{x \rightarrow x_0} (f(x) - f(x_0)) = \lim_{x \rightarrow x_0} \left[ \frac{f(x) - f(x_0)}{x - x_0} \right](x - x_0) = f'(x_0) \cdot 0 = 0 \nonumber
\end{align}
\end{proof}
\end{theorem}

\begin{definition} \hlt{Product Formula}\\
The derivative of product of two functions.
\begin{align}
(u \cdot v)' = u' \cdot v + u \cdot v' \nonumber
\end{align}
\end{definition}

\begin{definition} \hlt{Quotient Formula}\\
The derivate of ratio of two differentiable functions. Let $v(x) \neq 0$, then
\begin{align}
\left( \frac{u}{v} \right)' = \frac{u'v - uv'}{v^2} \nonumber
\end{align}
\end{definition}

\begin{definition} \hlt{Chain Rule}\\
The derivative of composition of two functions.
\begin{align}
\frac{d}{dt} f(g(t)) = f'(g(t))g'(t) \nonumber
\end{align}
May be used to differentiate implicitly defined functions.
\end{definition}

\begin{theorem} \hlt{Taylor's Theorem}\\
Let $k \geq 1$ be an integer and let function $f: \R \rightarrow \R$ be $k$-times differentiable at point $a \in \R$. Then there exists a function $h_k: \R \rightarrow \R$ such that
\begin{align}
f(x) &= \sum\limits_{i=0}^k \frac{f^{(i)}(a)}{i!} (x-a)^i + h_k (x) (x-a)^k \nonumber
\end{align}
where $\lim\limits_{x \rightarrow a} h_k(x) = 0$ is the Peano form of the remainder.
\end{theorem}

\begin{theorem} \hlt{Mean Value Theorem}\\
Let $f: [a,b] \rightarrow \R$ be a continuous function on closed interval $[a,b]$ and differentiable on open interval $(a,b)$ where $a < b$. Then $\exists c \in (a,b)$ such that
\begin{align}
f'(c) = \frac{f(b) - f(a)}{b-a} \nonumber
\end{align}
\end{theorem}

\begin{theorem} \hlt{First Fundamental Theorem of Calculus}\\
Let $f$ be continuous real-valued function on closed interval $[a,b]$ Let $F$ be function defined for all $x \in [a,b]$, by
\begin{align}
F(x) = \int_{a}^x f(t) dt \nonumber
\end{align}
Then $F$ is uniformly continuous on $[a,b]$ and differentiable on open interval $(a,b)$, and
\begin{align}
F'(x) = f(x) \nonumber
\end{align}
for all $x$ in $(a,b)$ so $F$ is an antiderivative of $f$.
\end{theorem}

\begin{remark} \hlt{Change of Variable}\\
If $f(x) = g(u(x))$, then $du = u'(x) dx$, hence
\begin{align}
\int g(u) du = \int g(u(x)) u'(x) dx = \int f(x) u'(x) dx \nonumber
\end{align}
For definite integrals,
\begin{align}
\int_{x_1}^{x_2} f(x) u'(x) dx = \int_{u_1}^{u_2} g(u) d(u) \nonumber
\end{align}
where $u_1 = u(x_1), u_2 = u(x_2)$.
\end{remark}

\begin{theorem} \hlt{Second Fundamental Theorem of Calculus}\\
Let $f: \R \rightarrow \R$ on $[a,b]$ and $F$ a continuous function on $[a,b]$ which is an antiderivative of $f$ in $(a,b)$,
\begin{align}
F'(x) = f(x) \nonumber
\end{align}
If $f$ is Riemann integrable on $[a,b]$, then
\begin{align}
\int_{a}^b f(x) dx = F(b) - F(a) \nonumber
\end{align}
\end{theorem}

\begin{remark} \hlt{Integration by Trigonometric Substitution}\\
Evaluate trigonometric integral, work backwards to original variable by using $\text{trig}(\text{trig}^{-1})$ with right-triangle.
\begin{enumerate}[label=\roman*.]
\setlength{\itemsep}{0pt}
\item $a^2 - x^2$ use $x = a \sin u$ and $dx = a \cos u \ du$
\item $a^2 + x^2$ use $x = a \tan u$ and $dx = a \sec^2 u \ du$
\item $x^2 - a^2$ use $x = a \sec u$ and $dx = a \sec u \tan u \ du$
\end{enumerate}
\end{remark}

\begin{remark} \hlt{Systematic Integration of Rational Functions}\\
For a rational function $P/Q$ (ratio of polynomials), by method of partial fractions, factor $Q$ completely into linear and quadratic factors, and write $P/Q$ as a sum of simpler terms.
\begin{enumerate}[label=\roman*.]
\setlength{\itemsep}{0pt}
\item If $Q(x) = (x-a)^m (x-b)^n$, then the expression is
\begin{align}
\frac{A_1}{x-a} + \frac{A_2}{(x-a)^2} + \cdots + \frac{A_m}{(x-a)^m} + \frac{B_1}{x-b} + \frac{B_2}{(x-b)^2} + \cdots + \frac{B_n}{(x-b)^n} \nonumber
\end{align}
\item If there are quadratic factor such as $(Ax^2 + Bx + C)^p$, then there are terms
\begin{align}
\frac{a_1 x + b_1}{Ax^2 + Bx + C} + \frac{a_2 x + b_2}{(Ax^2 + Bx + C)^2} + \cdots + \frac{a_p x + b_p}{(Ax^2 + Bx + C)^p} \nonumber
\end{align}
for each factor. To integrate, complete the square and make trigonometric substitution.
\end{enumerate}
\end{remark}

\begin{remark} \hlt{Integration by Parts}
\begin{align}
\int_b^a uv' \ dx = uv\vert_a^b - \int_a^b u'v \ dx \nonumber
\end{align}
This is used when $u'v$ is simpler than $uv'$ (often the case when $u'$ is simpler than $u$).
\end{remark}

\begin{remark} \hlt{Arc-Length}\\
Given $ds = \sqrt{dx^2 + dy^2}$, depending on whether to integrate with respect to $x,t,y$, this is
\begin{align}
ds = \sqrt{1+ \left(\frac{dy}{dx}\right)^2} \ dx; \ \ \ ds = \sqrt{\left(\frac{dx}{dt}\right)^2 + \left(\frac{dy}{dt}\right)^2} \ dt; \ \ \ ds = \sqrt{\left(\frac{dx}{dy}\right)^2 + 1} \ dy \nonumber
\end{align}
\end{remark}

\begin{remark} \hlt{Surface Area for Surface of Revolution}
\begin{enumerate}[label=\roman*.]
\setlength{\itemsep}{0pt}
\item Around $x$-axis: (requires formula for $y = y(x)$)
\begin{align}
2\pi y \ ds = 2 \pi y \sqrt{1+ \left(\frac{dy}{dx}\right)^2} \ dx \nonumber
\end{align}
\item Around $y$-axis: (requires formula for $x = x(y)$)
\begin{align}
2\pi x \ ds = 2 \pi x \sqrt{\left(\frac{dx}{dy}\right)^2 + 1} \ dy \nonumber
\end{align}
\end{enumerate}
\end{remark}

\begin{remark} \hlt{Polar Coordinates}\\
Given $x = r \cos \theta$, $y = r \sin \theta$ (or more rarely $r = \sqrt{x^2 + y^2} \theta$, $\theta = \tan^{-1}(y/x)$)
\begin{enumerate}[label=\roman*.]
\setlength{\itemsep}{0pt}
\item Find the polar equation for a curve from its equation in $(x,y)$ variables by substitution.
\item Sketch curves given in polar coordinates and find range of $\theta$ for integration.
\item Integration area formula is then
\begin{align}
\int_{\theta_1}^{\theta_2} \frac{1}{2} r^2 \ d \theta \nonumber
\end{align}
\end{enumerate}
\end{remark}

\begin{theorem} \hlt{L'Hopital Rule} \\
For functions $f$ and $g$ defined on open interval $I$ and differentiable on $I \backslash \{c\}$ for a accumulation point $c$ of $I$, if
\begin{align}
\lim\limits_{x \rightarrow c} f(x) = \lim\limits_{x \rightarrow c} g(x) = 0 \ \ \ \text{or} \ \ \ \pm \infty \nonumber
\end{align}
and $g'(x) \neq 0$ for all $x \in I \backslash \{ c \}$, and $\lim\limits_{x \rightarrow c} \frac{f'(x)}{g'(x)}$ exists, then
\begin{align}
\lim\limits_{x \rightarrow c} \frac{f(x)}{g(x)} = \lim\limits_{x \rightarrow c} \frac{f'(x)}{g'(x)} \nonumber
\end{align}
\end{theorem}